\documentclass[10pt,a4paper]{amsart}
\usepackage[margin=19mm]{geometry}
\usepackage{amsmath,amssymb,mathtools}
\usepackage[scr=rsfs,cal=euler]{mathalfa}
\usepackage{xfrac}
\usepackage[unicode,hidelinks,bookmarks=false]{hyperref}
\newcommand{\ie}{i.e.\ }
\newcommand{\cf}{cf.\ }
\newcommand{\resp}{respectively\ }
\newcommand{\Subsection}{\subsection}
\providecommand{\nobreakdash}{\nobreak\hbox{-}}
\setlength{\emergencystretch}{2em}
\multlinegap0pt
\newtheorem{proposition}{Proposition}
\title{Anomaly Detection from a Tensor Train Perspective}
\author{Alejandro Mata Ali, Aitor Moreno Fdez. de Leceta, Jorge L\'opez Rubio}
\date{Source: arXiv:2409.15030v2 (CC BY 4.0)}
\begin{document}
\maketitle

\noindent\emph{Source and licence.} Anomaly Detection from a Tensor Train Perspective. Version arXiv:2409.15030v2 (CC BY 4.0), distributed by arXiv under the Creative Commons Attribution 4.0 International licence, http://creativecommons.org/licenses/by/4.0/, as recorded in the arXiv licence metadata for this exact version and shown on the abstract page https://arxiv.org/abs/2409.15030v2. Full licence notice: \texttt{licenses/arxiv-2409.15030-CC-BY-4.0.txt}.\par\smallskip

\noindent\textit{Editorial scope.} Counted unit: the article's proposition proposition (source0.tex lines 263--278). The article's complete original proof of it is delivered with it (source0.tex lines 280--298, one \texttt{proof} environment of 19 lines). No mathematics is written, repaired or re-derived: the statement and the proof are the article's own lines, in the article's own notation, and no retained line is substituted or re-flowed.  No further source-local context is needed: every cross-reference of every retained line resolves inside the two delivered spans.  Nothing was omitted from any retained span, so the package carries no omission record.  No retained line contains a citation, so the wrapper carries no bibliography and no bibliography entry.  No retained line contains a figure environment, an image-inclusion command or drawing code, so no image is delivered and the package declares no asset dependency.  Editorial formatting only, in addition: an AMS \texttt{amsart} preamble that restates the article's own declarations and macros used by the retained lines, namely the environments \texttt{proposition} on separate counters, together with the standard packages those lines need.  The retained environments print in this document as Proposition 1 in order of appearance, whereas the article prints the same items with the numbering of its own full document.  The complete native source of the article as distributed in the arXiv e-print for this exact version is bundled unchanged as \texttt{sources/165673/source0.tex}.
\begin{proposition}
Let $Q\in\mathbb{R}^{D\times r}$ satisfy $Q^\top Q=I_r$. For any
$Y\in\mathbb{R}^{D\times m}$, the minimizer of
\[
    \min_{H\in\mathbb{R}^{r\times m}}
    \lVert Y-QH\rVert_F^2
\]
is
\[
    H^\star = Q^\top Y.
\]
Moreover, the minimum value is
\[
    \lVert (I-QQ^\top)Y\rVert_F^2.
\]
\end{proposition}

\begin{proof}
Since $Q^\top Q=I_r$, we may decompose
\[
    Y-QH
    =
    Q(Q^\top Y-H) + (I-QQ^\top)Y.
\]
The two terms on the right-hand side are orthogonal with respect to the
Frobenius inner product. Hence
\[
    \lVert Y-QH\rVert_F^2
    =
    \lVert Q^\top Y-H\rVert_F^2
    +
    \lVert (I-QQ^\top)Y\rVert_F^2.
\]
The second term does not depend on $H$, and the first term is minimized
uniquely by $H=Q^\top Y$.
\end{proof}

\end{document}
